% =================================================
% Содержание темы
% Подключается и к цветной, и к чёрно-белой версии.
% =================================================

\begin{center}
  {\fontsize{25}{30}\selectfont\bfseries\sffamily\color{Primary}
    Рациональное уравнение с модулем и параметром}

  \vspace{0.3em}

  {\large\sffamily\color{GrayText}
    Три метода решения одной задачи}
\end{center}

\vspace{0.45em}

\section{Постановка задачи}

\begin{problembox}
  Найдите все значения параметра \(a\), при каждом из которых уравнение
  \[
    \frac{|3x|-2x-2-a}{x^2-2x-a}=0
  \]
  имеет ровно два различных корня.
\end{problembox}

В этой задаче недостаточно исследовать только числитель. Найденный ноль
числителя не является корнем исходного уравнения, если при том же значении
\(x\) знаменатель обращается в ноль.

\begin{propertybox}
  Дробь равна нулю тогда и только тогда, когда при одном и том же значении
  переменной её числитель равен нулю, а знаменатель отличен от нуля:
  \[
    \frac{A(x)}{B(x)}=0
    \quad\Longleftrightarrow\quad
    \begin{cases}
      A(x)=0,\\
      B(x)\ne 0.
    \end{cases}
  \]
\end{propertybox}

Поэтому исходную задачу удобно разделить на два исследования.

\parameterlogicfigure

\Needspace{12\baselineskip}
\section{Подготовительные задания}

Перед основной задачей полезно отдельно отработать три приёма: ноль дроби,
раскрытие модуля и поиск общего корня двух выражений.

\subsection{Ноль рационального выражения}

\begin{trainingtask}[A. Два коротких уравнения]
  Решите уравнения:
  \begin{enumerate}
    \item \(\displaystyle \frac{x-3}{x+1}=0;\)
    \item \(\displaystyle \frac{(x-2)(x+3)}{x-2}=0.\)
  \end{enumerate}
  Во втором уравнении объясните, почему один из нулей числителя не является
  корнем исходной дроби.
\end{trainingtask}

\begin{notebox}[Проверка]
  В первом уравнении \(x=3\). Во втором при \(x=2\) знаменатель обращается
  в ноль, поэтому это значение не является корнем исходного уравнения.
  Единственный корень равен \(x=-3\).
\end{notebox}

\subsection{Модуль и количество корней}

\begin{trainingtask}[B. От отдельного уравнения к исследованию]
  \begin{enumerate}
    \item Решите уравнение \(3|x|-2x=5\).
    \item Определите количество корней уравнения
      \[
        3|x|-2x=b
      \]
      в зависимости от значения \(b\).
  \end{enumerate}
\end{trainingtask}

Так как
\[
  3|x|-2x=
  \begin{cases}
    -5x, & x<0,\\
    x, & x\ge 0,
  \end{cases}
\]
обе ветви принимают только неотрицательные значения. Отсюда получаем:

\begin{center}
  \renewcommand{\arraystretch}{1.28}
  \begin{tabular}{cc}
    \toprule
    \rowcolor{TableHeader}
    Значение \(b\) & Количество корней \\
    \midrule
    \(b<0\) & \(0\) \\
    \(b=0\) & \(1\) \\
    \(b>0\) & \(2\) \\
    \bottomrule
  \end{tabular}
\end{center}

В первом пункте \(b=5>0\), поэтому получаются два корня:
\[
  x=5,
  \qquad
  x=-1.
\]

\subsection{Общий корень}

\begin{trainingtask}[C. Подстановка корня]
  При каких значениях \(a\) число \(x=a+1\) является корнем уравнения
  \[
    x^2-x-a=0?
  \]
\end{trainingtask}

Подставим \(x=a+1\):
\[
  (a+1)^2-(a+1)-a=a^2.
\]
Следовательно, общий корень возникает только при \(a=0\).

\begin{conclusionbox}[Что переносим в основную задачу]
  Сначала определяем, когда числитель имеет нужное количество нулей. Затем
  подстановкой или решением системы находим параметры, при которых эти нули
  одновременно обращают знаменатель в ноль.
\end{conclusionbox}

\Needspace{14\baselineskip}
\section{Метод 1. Прямое аналитическое решение}

\begin{methodintro}
  Раскрываем модуль по промежуткам, находим два возможных корня числителя,
  а затем непосредственно проверяем знаменатель для каждого корня.
\end{methodintro}

Дробь равна нулю тогда и только тогда, когда
\[
  \begin{cases}
    |3x|-2x-2-a=0,\\
    x^2-2x-a\ne 0.
  \end{cases}
\]

\stepheading{Шаг 1}{Исследуем числитель}

При \(x\ge 0\) решаем систему, сохраняя условие выбранной ветви модуля:
\[
  \begin{cases}
    3x-2x-2-a=0,\\
    x\ge 0
  \end{cases}
  \quad\Longleftrightarrow\quad
  \begin{cases}
    x_1=a+2,\\
    a\ge -2.
  \end{cases}
\]

При \(x<0\) аналогично получаем:
\[
  \begin{cases}
    -3x-2x-2-a=0,\\
    x<0
  \end{cases}
  \quad\Longleftrightarrow\quad
  \begin{cases}
    x_2=-\dfrac{a+2}{5},\\
    a>-2.
  \end{cases}
\]

\begin{conclusionbox}
  Числитель имеет два различных нуля тогда и только тогда, когда
  \[
    \boxed{a>-2}.
  \]
\end{conclusionbox}

\begin{notebox}
  При \(a=-2\) обе формулы дают число \(0\), но отрицательная ветвь требует
  \(x<0\). Поэтому в этом случае числитель имеет только один ноль.
\end{notebox}

\stepheading{Шаг 2}{Проверяем корень \(x_1=a+2\)}

\begin{align*}
  x_1^2-2x_1-a
  &=(a+2)^2-2(a+2)-a\\
  &=a^2+a\\
  &=a(a+1).
\end{align*}

Следовательно, при \(a=-1\) или \(a=0\) значение
\(x_1=a+2\) обращает знаменатель в ноль и потому не является корнем
исходного уравнения.

\stepheading{Шаг 3}{Проверяем корень \(x_2=-\dfrac{a+2}{5}\)}

\begin{align*}
  x_2^2-2x_2-a
  &=\frac{(a+2)^2}{25}+\frac{2(a+2)}{5}-a\\
  &=\frac{a^2-11a+24}{25}\\
  &=\frac{(a-3)(a-8)}{25}.
\end{align*}

Следовательно, при \(a=3\) или \(a=8\) значение
\(x_2=-\dfrac{a+2}{5}\) обращает знаменатель в ноль и потому не является
корнем исходного уравнения.

\begin{conclusionbox}[Проверка знаменателя]
  Оба нуля числителя допустимы, если
  \[
    a\notin\{-1,0,3,8\}.
  \]
\end{conclusionbox}

\begin{resultbox}
  Объединяем условие \(a>-2\) с ограничениями знаменателя:
  \[
    \boxed{a\in(-2,+\infty)\setminus\{-1,0,3,8\}}.
  \]
\end{resultbox}

\Needspace{14\baselineskip}
\section{Метод 2. Графико-аналитическое решение}

\begin{methodintro}
  Количество нулей числителя определяем по графику. Допустимость найденных
  корней по-прежнему проверяем аналитически.
\end{methodintro}

Уравнение числителя перепишем в виде
\[
  3|x|-2x=a+2.
\]
Обозначим \(b=a+2\). При изменении параметра \(a\) на ту же величину
изменяется и \(b\). Поэтому, двигая горизонтальную прямую \(y=b\) вверх
или вниз, мы фактически изменяем параметр \(a\): чем больше \(a\), тем
выше расположена прямая.

Рассмотрим график функции
\[
  y=3|x|-2x=
  \begin{cases}
    -5x, & x<0,\\
    x, & x\ge 0.
  \end{cases}
\]
Количество корней равно количеству пересечений этого графика с прямой \(y=b\).

\parametergraphfigure

По графику видно:
\begin{itemize}
  \item \(\boldsymbol{a<-2}\): прямая проходит ниже точки излома
    \(O(0;0)\), точек пересечения нет;
  \item \(\boldsymbol{a=-2}\): прямая проходит через точку излома,
    имеется одна точка пересечения;
  \item \(\boldsymbol{a>-2}\): прямая расположена выше точки излома и
    пересекает обе ветви, имеются две точки пересечения.
\end{itemize}

\begin{conclusionbox}
  Условие
  \[
    \boxed{a>-2}
  \]
  необходимо и достаточно для того, чтобы числитель имел два различных
  нуля. Для исходного рационального уравнения оно пока является только
  необходимым: ещё нужно проверить знаменатель.

  Абсциссы точек пересечения равны
  \[
    x_1=a+2,
    \qquad
    x_2=-\frac{a+2}{5}.
  \]
\end{conclusionbox}

Остаётся проверить знаменатель. Результаты удобно собрать в таблицу.

\begin{center}
  \renewcommand{\arraystretch}{1.42}
  \begin{tabularx}{\textwidth}{
      >{\centering\arraybackslash}p{0.22\textwidth}
      >{\centering\arraybackslash}p{0.36\textwidth}
      >{\centering\arraybackslash}X
    }
    \toprule
    \rowcolor{TableHeader}
    Ноль числителя & Знаменатель после подстановки &
      Знаменатель равен нулю при \\
    \midrule
    \(x_1=a+2\) & \(a(a+1)\) & \(-1,\ 0\) \\
    \addlinespace[0.25em]
    \(x_2=-\dfrac{a+2}{5}\) & \(\dfrac{(a-3)(a-8)}{25}\) & \(3,\ 8\) \\
    \bottomrule
  \end{tabularx}
\end{center}

\begin{warningbox}
  График отвечает на вопрос о количестве нулей \emph{числителя}. Область
  определения дроби на нём не отражена, поэтому проверка знаменателя остаётся
  обязательной.
\end{warningbox}

\begin{resultbox}
  \[
    \boxed{a\in(-2,+\infty)\setminus\{-1,0,3,8\}}.
  \]
\end{resultbox}

\Needspace{14\baselineskip}
\section{Метод 3. Исключение общих корней}

\begin{methodintro}
  Как было показано ранее, при \(a>-2\) числитель имеет ровно два различных
  нуля. Теперь найдём значения параметра, при которых хотя бы один из них
  обращает знаменатель в ноль, то есть числитель и знаменатель имеют общий
  корень.
\end{methodintro}

Рассмотрим систему
\[
  \begin{cases}
    |3x|-2x-2-a=0,\\
    x^2-2x-a=0.
  \end{cases}
\]
Из первого уравнения \(a=3|x|-2x-2\), а из второго \(a=x^2-2x\). Поэтому
\[
  \begin{aligned}
    3|x|-2x-2&=x^2-2x,\\
    3|x|-2&=x^2.
  \end{aligned}
\]

\begin{minipage}[t]{0.485\textwidth}
  \begin{casebox}{Случай \(x\ge 0\)}
    Здесь \(|x|=x\), поэтому
    \[
      \begin{aligned}
        x^2-3x+2&=0,\\
        (x-1)(x-2)&=0.
      \end{aligned}
    \]
    Получаем \(x=1\) или \(x=2\). Тогда
    \[
      \begin{aligned}
        x=1&\Rightarrow a=-1,\\
        x=2&\Rightarrow a=0.
      \end{aligned}
    \]
  \end{casebox}
\end{minipage}\hfill
\begin{minipage}[t]{0.485\textwidth}
  \begin{casebox}{Случай \(x<0\)}
    Здесь \(|x|=-x\), поэтому
    \[
      \begin{aligned}
        x^2+3x+2&=0,\\
        (x+1)(x+2)&=0.
      \end{aligned}
    \]
    Получаем \(x=-1\) или \(x=-2\). Тогда
    \[
      \begin{aligned}
        x=-1&\Rightarrow a=3,\\
        x=-2&\Rightarrow a=8.
      \end{aligned}
    \]
  \end{casebox}
\end{minipage}

\begin{conclusionbox}[Общие корни и окончательный ответ]
  Числитель и знаменатель имеют общий корень только при
  \[
    a\in\{-1,0,3,8\}.
  \]
  Исключаем эти значения из условия \(a>-2\):
  \[
    \boxed{a\in(-2,+\infty)\setminus\{-1,0,3,8\}}.
  \]
\end{conclusionbox}

\Needspace{14\baselineskip}
\section{Итог и сравнение методов}

Ответ удобно представить на числовой прямой.

\parameternumberlinefigure

\begin{center}
  \renewcommand{\arraystretch}{1.38}
  \begin{tabularx}{\textwidth}{
      >{\raggedright\arraybackslash}p{0.24\textwidth}
      >{\raggedright\arraybackslash}p{0.37\textwidth}
      >{\raggedright\arraybackslash}X
    }
    \toprule
    \rowcolor{TableHeader}
    Метод & Главная идея & Что особенно хорошо показывает \\
    \midrule
    Аналитический
      & Раскрыть модуль и проверить каждый корень подстановкой
      & Полную последовательность вычислений \\
    \addlinespace[0.25em]
    Графико-аналитический
      & Найти количество пересечений с горизонтальной прямой
      & Почему при \(a>-2\) числитель имеет два различных нуля \\
    \addlinespace[0.25em]
    Общие корни
      & Решить систему «числитель равен нулю, знаменатель равен нулю»
      & Откуда появляются четыре исключаемых значения \\
    \bottomrule
  \end{tabularx}
\end{center}

\begin{propertybox}[Главная идея задачи]
  Исследование рационального уравнения с параметром естественно делится на
  два слоя:
  \[
    \text{количество нулей числителя}
    \quad+\quad
    \text{допустимость каждого нуля}.
  \]
\end{propertybox}

\subsection{Полная классификация количества корней}

\begin{center}
  \renewcommand{\arraystretch}{1.38}
  \begin{tabularx}{\textwidth}{
      >{\centering\arraybackslash}p{0.48\textwidth}
      >{\centering\arraybackslash}X
    }
    \toprule
    \rowcolor{TableHeader}
    Значение параметра & Количество корней \\
    \midrule
    \(a<-2\) & \(0\) \\
    \(a=-2\) & \(1\) \\
    \(a\in\{-1,0,3,8\}\) & \(1\) \\
    \(a>-2,\ a\notin\{-1,0,3,8\}\) & \(2\) \\
    \bottomrule
  \end{tabularx}
\end{center}

\Needspace{12\baselineskip}
\section{Задания для самостоятельной работы}

\begin{practicebox}[Задание 1. Близкая модель]
  Найдите все значения \(a\), при которых уравнение
  \[
    \frac{|x|-a}{x-2}=0
  \]
  имеет ровно два различных корня.
\end{practicebox}

\begin{practicebox}[Задание 2. Проверка знаменателя]
  Найдите все значения \(a\), при которых уравнение
  \[
    \frac{2|x|-x-a}{x^2-a}=0
  \]
  имеет ровно два различных корня.
\end{practicebox}

\begin{practicebox}[Задание 3. Аналог основной задачи]
  Найдите все значения \(a\), при которых уравнение
  \[
    \frac{|2x|-x-1-a}{x^2-x-a}=0
  \]
  имеет ровно два различных корня.
\end{practicebox}

\begin{practicebox}[Задание 4. Анализ ошибки]
  Ученик исследовал только числитель исходного уравнения, получил условие
  \(a>-2\) и сразу записал его в ответ. Объясните ошибку и завершите решение.
\end{practicebox}


\begin{practicebox}[Задание 5*. Параметр внутри модуля]
  Найдите все значения \(a\), при которых уравнение
  \[
    \frac{2|x-a|-x-2}{x^2-1}=0
  \]
  имеет ровно два различных корня.

  Здесь при изменении параметра график \(y=2|x-a|\) перемещается вдоль оси
  \(Ox\), поэтому готовый шаблон основной задачи нельзя применять
  механически.
\end{practicebox}

\begin{questionsbox}
  \begin{enumerate}
    \item Почему значение \(a=-2\) не входит в ответ основной задачи?
    \item Что происходит с числом корней при \(a=-1\), \(0\), \(3\) и \(8\)?
    \item Как по графику функции \(y=3|x|-2x\) увидеть условие \(a>-2\)?
    \item Почему третий метод сразу находит именно исключаемые значения параметра?
    \item Какой из трёх методов удобнее использовать для проверки вычислений?
  \end{enumerate}
\end{questionsbox}

\begin{notebox}[Ответы к самостоятельной работе]
  \begin{enumerate}
    \item \(a\in(0,+\infty)\setminus\{2\}\).
    \item \(a\in(0,+\infty)\setminus\{1,9\}\).
    \item \(a\in(-1,+\infty)\setminus\{0,2\}\).
    \item Нужно дополнительно исключить \(-1\), \(0\), \(3\) и \(8\).
    \item \(a\in(-2,+\infty)\setminus
      \left\{-\dfrac32,-\dfrac12,\dfrac52\right\}\).
  \end{enumerate}
\end{notebox}
